\documentclass[11pt]{article}

\usepackage[margin=0.75in]{geometry}
\usepackage[hidelinks]{hyperref}
\usepackage{enumitem}
\usepackage{titlesec}
\usepackage{parskip}

\pagestyle{plain}
\setlength{\parindent}{0pt}
\emergencystretch=1.5em

\titleformat{\section}{\large\bfseries\scshape}{}{0em}{}[\titlerule]
\titlespacing{\section}{0pt}{10pt}{6pt}

\newcommand{\entry}[2]{\noindent\textbf{#1}\hfill #2\par}
\newcommand{\subentry}[2]{\noindent #1\hfill #2\par}

\begin{document}

\begin{center}
    {\LARGE\bfseries Longtian Bao}\\[4pt]
    \href{mailto:lobao@uchicago.edu}{lobao@uchicago.edu} \,|\,
    \href{https://github.com/LongtianBao}{github.com/LongtianBao} \,|\,
    \href{https://scholar.google.com/citations?user=xX-1jwkAAAAJ&hl=en}{Google Scholar}
\end{center}

\section{Education}
\entry{University of Chicago}{Chicago, IL}
\subentry{Ph.D. in Computer Science}{2026 -- Present}
\entry{University of California, San Diego}{La Jolla, CA}
\subentry{M.S. in Computer Science (GPA: 4.0/4.0)}{2024 -- 2026}
\subentry{B.S. in Computer Engineering (GPA: 3.8/4.0)}{2020 -- 2024}

\section{Publications}
{\small * denotes equal contribution.}

\begin{enumerate}[label={[\arabic*]}, leftmargin=2.2em, itemsep=4pt]
    \item \textbf{Longtian Bao}*, Jianyou Wang*, Yang Zhang*, Youze Zheng*, Ramamohan Paturi.
    ``Question Begets Question: Self-Evolving Curriculum for Reinforcement Fine-Tuning on Competition Mathematics.''
    \textit{arXiv preprint arXiv:2608.01522}, 2026.
    \href{https://arxiv.org/abs/2608.01522}{[Paper]}

    \item \textbf{Longtian Bao}, Jiajun Xi, Jingyi Zhang, Zhengan Cheng, Junfei Liu, Jinya Jiang.
    ``Multi-Agent LLM Serving Is a Full-Stack Problem: A Small Step Towards Algorithm--System Co-Design.''
    \textit{Under review}, 2026.

    \item Youze Zheng, Jianyou Wang, Yuhan Chen, Matthew Feng, \textbf{Longtian Bao}, Hanyuan Zhang, Maxim Khan, Aditya K. Sehgal, Christopher D. Rosin, Umber Dube, Ramamohan Paturi.
    ``DeepImagine: Learning Biomedical Reasoning via Successive Counterfactual Imagining.''
    \textit{arXiv preprint arXiv:2604.23054}, 2026.
    \href{https://arxiv.org/abs/2604.23054}{[Paper]}

    \item Jianyou Wang, Youze Zheng, \textbf{Longtian Bao}, Hanyuan Zhang, Qirui Zheng, Yuhan Chen, Yang Zhang, Matthew Feng, Maxim Khan, Aditya K. Sehgal, Christopher D. Rosin, Ramamohan Paturi, Umber Dube, Leon Bergen.
    ``CT Open: An Open-Access, Uncontaminated, Live Platform for the Open Challenge of Clinical Trial Outcome Prediction.''
    \textit{arXiv preprint arXiv:2604.16742}, 2026.
    \href{https://arxiv.org/abs/2604.16742}{[Paper]}

    \item Weili Cao*, Jianyou Wang, Youze Zheng*, \textbf{Longtian Bao}*, Qirui Zheng, Taylor Berg-Kirkpatrick, Ramamohan Paturi, Leon Bergen.
    ``Single-Pass Document Scanning for Question Answering.''
    \textit{Conference on Language Modeling (COLM)}, 2025. \textbf{(Oral Spotlight, Top 2\%)}.
    \href{https://arxiv.org/abs/2504.03101}{[Paper]}
    \href{https://github.com/weilicao/SPScanner}{[Code]}

    \item Jianyou Wang*, Weili Cao*, \textbf{Longtian Bao}, Youze Zheng, Gil Pasternak, Kaicheng Wang, Xiaoyue Wang, Ramamohan Paturi, Leon Bergen.
    ``Measuring Risk of Bias in Biomedical Reports: The RoBBR Benchmark.''
    \textit{Empirical Methods in Natural Language Processing (EMNLP)}, 2025.
    \href{https://arxiv.org/abs/2411.18831}{[Paper]}
    \href{https://github.com/RoBBR-Benchmark/RoBBR}{[Code]}
\end{enumerate}

\section{Research Experience}
\entry{Research Assistant --- Laboratory for Emerging Intelligence, UC San Diego}{La Jolla, CA}
\subentry{Advisors: Prof. Leon Bergen and Prof. Ramamohan Paturi}{Apr. 2024 -- Present}
\begin{itemize}[leftmargin=1.2em, itemsep=2pt, topsep=4pt]
    \item Re-architected a Mamba-2 State-Space Model (SSM) by replacing its language modeling head with a classification head, creating a subquadratic model that preserves global context for efficient sentence retrieval in documents as long as 256k tokens [5].
    \item Designed and implemented a novel link-based synthetic data generation pipeline where an LLM identifies and leverages natural, long-range dependencies within documents to create coherent training examples for long-context question answering [5].
    \item Developed an LLM-assisted annotation-validation pipeline to construct a new information retrieval benchmark from 41 existing QA datasets, enabling a fine-grained evaluation of model performance in retrieving relevant sentences from long documents [5].
    \item Developed the RoBBR benchmark, a new resource for assessing risk-of-bias in biomedical research, derived from expert judgments on over 500 scientific studies to evaluate LLM capabilities in analyzing methodological strength [6].
    \item Designed two novel subtasks---Support Sentence Retrieval (SSR) and Support Judgment Selection (SJS)---to disentangle and independently measure the information retrieval and reasoning abilities of language models in the complex domain of bias assessment [6].
    \item Engineered an automatic, human-validated annotation pipeline using GPT-4 to create the ground truth for the retrieval task by decomposing expert judgments into fine-grained ``aspects'' and mapping them precisely to source sentences in biomedical papers [6].
\end{itemize}

\section{Teaching Experience}
\entry{Instructional Apprentice (Tutor) --- UC San Diego}{La Jolla, CA}
\subentry{Selected Quarters:}{Spring 2021 -- Winter 2024}
\begin{itemize}[leftmargin=1.2em, itemsep=2pt, topsep=4pt]
    \item \textbf{CSE 12 - Data Structures:} Spring 2021, Winter 2023, Winter 2024
    \item \textbf{CSE 11 - Accelerated Introduction to Programming:} Winter 2022, Fall 2022, Spring 2023
    \item \textbf{CSE 30 - Computer Organization:} Fall 2023
    \item \textbf{CSE 15L - Software Tools \& Techniques Lab:} Spring 2021
\end{itemize}

\section{Projects}
\entry{UCSD Autograder Platform}{Fall 2023 -- Spring 2024}
\begin{itemize}[leftmargin=1.2em, itemsep=2pt, topsep=4pt]
    \item Engineered and maintained a full-stack educational platform used by over 1,000 students across multiple departments, managing a database with millions of rows of student data.
    \item Developed key features including a containerized autograder for homework submissions and a real-time, ticket-based queue system for teaching assistant office hours.
\end{itemize}

\entry{Autonomous AI Racing Car}{Winter 2024}
\begin{itemize}[leftmargin=1.2em, itemsep=2pt, topsep=4pt]
    \item Built and configured a 1/10 scale autonomous vehicle, integrating a Jetson Nano single-board computer with a VESC, LIDAR, and OAK-D AI camera.
    \item Utilized the DonkeyCar framework and TensorFlow to train a behavioral cloning model, enabling the car to autonomously navigate a racetrack by collecting and learning from human driving data.
\end{itemize}

\entry{DaCe Framework Reproducibility --- Student Cluster Competition}{Fall 2022}
\begin{itemize}[leftmargin=1.2em, itemsep=2pt, topsep=4pt]
    \item Evaluated the reproducibility of the DaCe high-performance Python framework on NPBench benchmarks using a novel AMD-based hardware configuration.
    \item Ported and debugged several NumPy benchmarks to the DaCe framework, analyzing performance speedups and identifying portability limitations.
\end{itemize}

\section{Skills}
\textbf{Programming:} Python (advanced), C/C++, Java, JavaScript/TypeScript, Bash\\
\textbf{ML \& NLP/IR:} PyTorch, TensorFlow, Hugging Face Transformers, State-Space Models (Mamba-2), reinforcement fine-tuning (GRPO), long-context QA \& sentence retrieval, OpenCV\\
\textbf{Backend \& Data:} Flask/FastAPI, PostgreSQL, REST APIs, real-time queue systems, basic PySpark\\
\textbf{Systems \& DevOps:} Docker, Linux/SSH, Git/GitHub, AWS (EC2/S3), Distributed Training, Weights \& Biases\\
\textbf{Web/Frontend:} React\\
\textbf{Robotics/Edge:} NVIDIA Jetson, LIDAR, OAK-D, VESC, DonkeyCar framework\\
\textbf{LLM Tooling:} OpenAI/Anthropic/Google APIs

\end{document}
